% ============================================================
%  R-Matrix Derivation: IMU Quaternion → World-Frame Rotation
%  realtime-vision-control / sagittal plane projection
%  Author: (fill in)   Date: 2026-04-21
% ============================================================
\documentclass[11pt, a4paper]{article}

\usepackage{amsmath, amssymb, amsthm}
\usepackage{bm}          % \bm for bold math
\usepackage{geometry}
\usepackage{hyperref}
\usepackage{booktabs}
\usepackage{graphicx}
\usepackage{xcolor}
\usepackage{enumitem}

\geometry{margin=2.5cm}

\newcommand{\R}{\mathbb{R}}
\newcommand{\vv}[1]{\bm{#1}}          % vector: \vv{p}
\newcommand{\mat}[1]{\mathbf{#1}}     % matrix: \mat{R}
\newcommand{\quat}[1]{\bm{#1}}        % quaternion: \quat{q}
\newcommand{\frame}[1]{\mathcal{#1}}  % frame: \frame{C}
\newcommand{\norm}[1]{\left\lVert #1 \right\rVert}
\newcommand{\Rtrans}[2]{\mat{R}_{#1 \leftarrow #2}}  % R_{W←C}

\theoremstyle{definition}
\newtheorem{defn}{Definition}
\newtheorem{prop}{Proposition}
\newtheorem{remark}{Remark}

% ============================================================
\title{\textbf{World-Frame Rotation Matrix Derivation}\\
       \large IMU Quaternion $\to$ Gravity-Aligned Sagittal Projection}
\date{2026-04-21}
\begin{document}
\maketitle
\tableofcontents
\vspace{1em}

% ============================================================
\section{Coordinate Frames}
% ============================================================

We define three right-handed coordinate frames used throughout the system.

\begin{defn}[Camera Frame $\frame{C}$]
Origin at the ZED~X~Mini left optical centre.
\begin{align*}
  \hat{\vv{x}}_C &= \text{image right (horizontal)} \\
  \hat{\vv{y}}_C &= \text{image down  (vertical)} \\
  \hat{\vv{z}}_C &= \text{optical axis (into scene)}
\end{align*}
Units: metres. This is the standard OpenCV camera convention.
\end{defn}

\begin{defn}[World Frame $\frame{W}$]
Gravity-aligned reference frame, frozen after IMU warmup.
\begin{align*}
  \hat{\vv{X}}_W &= \text{subject's left}\ (+) \\
  \hat{\vv{Y}}_W &= \text{gravity direction (down)}\ (+) \\
  \hat{\vv{Z}}_W &= \text{horizontal forward (walking direction)}\ (+)
\end{align*}
$\hat{\vv{Y}}_W$ is exactly anti-parallel to the gravity reaction vector measured
by the IMU accelerometer at rest:
\[
  \vv{g}_W = \begin{pmatrix} 0 \\ +g \\ 0 \end{pmatrix}, \quad g \approx 9.81\;\text{m/s}^2
\]
\end{defn}

\begin{defn}[Body / Human Frame $\frame{B}$]
Patient-centred frame, computed once per standing calibration.
\begin{align*}
  \text{Origin}  &= \tfrac{1}{2}(\vv{p}^W_{\text{L-hip}} + \vv{p}^W_{\text{R-hip}})
                    \quad\text{(hip midpoint in world frame)} \\
  \hat{\vv{x}}_B &= \text{pelvis lateral, L-hip} \to \text{R-hip, } \perp \hat{\vv{Y}}_W \\
  \hat{\vv{y}}_B &= \hat{\vv{Y}}_W \quad\text{(gravity, unchanged)} \\
  \hat{\vv{z}}_B &= \hat{\vv{x}}_B \times \hat{\vv{y}}_B \quad\text{(patient forward)}
\end{align*}
The \emph{sagittal plane} is the $(\hat{\vv{y}}_B,\,\hat{\vv{z}}_B)$ plane passing through
the hip midpoint.
\end{defn}

% ============================================================
\section{Quaternion to Rotation Matrix}
% ============================================================

The ZED~SDK pose-tracking pipeline fuses IMU data and visual odometry to
produce a unit quaternion
\[
  \quat{q} = w + x\,\vv{i} + y\,\vv{j} + z\,\vv{k}, \qquad
  w^2 + x^2 + y^2 + z^2 = 1,
\]
representing the camera's orientation in the gravity-aligned world frame
(Hamilton convention, active rotation).

\begin{prop}[Rotation matrix from unit quaternion]
\[
  \Rtrans{W}{C}(\quat{q}) =
  \begin{pmatrix}
    1 - 2(y^2+z^2) & 2(xy - wz)     & 2(xz + wy) \\
    2(xy + wz)     & 1 - 2(x^2+z^2) & 2(yz - wx) \\
    2(xz - wy)     & 2(yz + wx)     & 1 - 2(x^2+y^2)
  \end{pmatrix}
\]
such that a point expressed in the camera frame transforms to the world frame by
\[
  \vv{p}^W = \Rtrans{W}{C}\,\vv{p}^C.
\]
\end{prop}

\begin{remark}
$\Rtrans{W}{C} \in SO(3)$ satisfies $\Rtrans{W}{C}^{\!\top}\Rtrans{W}{C} = \mat{I}$
and $\det(\Rtrans{W}{C}) = +1$.
The inverse (camera-from-world) is $\Rtrans{C}{W} = \Rtrans{W}{C}^{\!\top}$.
\end{remark}

% ============================================================
\section{Gravity Vector and Pitch / Roll Extraction}
% ============================================================

At rest, the IMU accelerometer measures the specific force
$\vv{a}_{\text{IMU}} = -\vv{g}_C$ (reaction opposing gravity in camera frame).
Hence the gravity direction in camera frame is
\[
  \hat{\vv{g}}_C
    = \frac{-\vv{a}_{\text{IMU}}}{\norm{\vv{a}_{\text{IMU}}}}
    = \Rtrans{C}{W}\,\hat{\vv{Y}}_W
    = \Rtrans{W}{C}^{\!\top}\begin{pmatrix}0\\1\\0\end{pmatrix}.
\]
This single vector constrains two degrees of freedom of $\Rtrans{W}{C}$
(pitch $\theta$ and roll $\phi$) but leaves yaw $\psi$ around $\hat{\vv{Y}}_W$
\emph{undetermined}.

\subsection*{Yaw Ambiguity and Resolution}
The ZED SDK resolves yaw from visual odometry during the warmup phase.
Indoors without a reliable magnetometer, the resulting yaw is the camera's
initial heading rather than geographic north --- which is acceptable here
because the sagittal plane only requires the pelvis-lateral direction, not
absolute compass bearing.

For the \emph{viewer calibration} (Section~\ref{sec:sagittal}), residual yaw
error is corrected by the patient-specific hip-vector measurement.

% ============================================================
\section{Averaging N Quaternion Samples (Warmup)}
\label{sec:avg_quat}
% ============================================================

During warmup we collect $N$ pose quaternions
$\{\quat{q}_1, \dots, \quat{q}_N\}$ while the camera is stationary.
Arithmetic mean followed by renormalisation fails when rotations span
a large arc (the mean may leave $S^3$). Instead we use the
\textbf{eigendecomposition estimator} \cite{Markley2007}:

\begin{prop}[Optimal quaternion average]
Form the $4 \times 4$ accumulator matrix
\[
  \mat{M} = \sum_{i=1}^{N} \quat{q}_i\,\quat{q}_i^{\!\top} \in \R^{4\times 4}.
\]
The maximum-likelihood average quaternion (under isotropic Gaussian noise on
$S^3$) is the unit eigenvector $\quat{q}^*$ of $\mat{M}$ corresponding to its
largest eigenvalue $\lambda_{\max}$.
\end{prop}

\noindent This is equivalent to finding the rotation closest (in Frobenius norm)
to all $N$ samples simultaneously, and always returns a valid unit quaternion.

\subsection*{Stillness Gate}
\label{sec:stillness}
Before accepting the average, we verify that the camera was genuinely stationary:
\[
  \theta_i = 2\arccos\!\bigl(\lvert \quat{q}_i \cdot \quat{q}^* \rvert\bigr)
  \quad \in [0°, 180°]
\]
is the geodesic angle on $S^3$ between sample $i$ and the mean.  The gate is
\[
  \max_{i}\,\theta_i \;<\; \theta_{\max} = 2°.
\]
If any sample exceeds $\theta_{\max}$, the buffer is discarded and warmup
restarts, prompting the operator to hold the camera still.

% ============================================================
\section{Static Rotation $\mat{R}$ and Its Maintenance}
% ============================================================

After successful warmup:
\[
  \boxed{
    \mat{R} \;=\; \Rtrans{W}{C}\!\bigl(\quat{q}^*\bigr)
    \quad\in\; SO(3), \quad \text{frozen for the session.}
  }
\]
All subsequent world-frame 3-D keypoint coordinates are computed by
\[
  \vv{p}^W = \mat{R}\,\vv{p}^C \in \R^3.
\]
This operation is performed on the GPU (\texttt{gpu\_postprocess.py},
\texttt{\_lift\_to\_3d\_v}) and adds \textbf{zero extra latency} beyond the
existing matrix-vector multiply already in the hot path.

\subsection*{Why a Static $\mat{R}$?}
Per-frame IMU updates would cost an extra GPU-to-CPU transfer and synchronisation
barrier, adding $\sim\!2$--$5$\,ms to a budget of $<20$\,ms end-to-end.
The camera is mounted on a rigid exoskeleton frame; in practice the camera does
not move relative to gravity during a walking trial.  A static $\mat{R}$ is
therefore both safe and necessary.

% ============================================================
\section{Sagittal Plane Projection}
\label{sec:sagittal}
% ============================================================

\subsection{Body-Frame Construction}

Let $\vv{h}_L^W,\,\vv{h}_R^W \in \R^3$ be the left and right hip keypoints in
the world frame.  Define the hip midpoint
\[
  \vv{o}^W = \tfrac{1}{2}\!\left(\vv{h}_L^W + \vv{h}_R^W\right)
\]
and the raw lateral axis
\[
  \tilde{\vv{x}}_B = \vv{h}_R^W - \vv{h}_L^W.
\]
Projecting out the gravity component gives a pelvis-lateral axis that is
guaranteed perpendicular to $\hat{\vv{Y}}_W$ even when the patient leans:
\[
  \hat{\vv{x}}_B
    = \frac{\tilde{\vv{x}}_B - (\tilde{\vv{x}}_B \cdot \hat{\vv{Y}}_W)\,\hat{\vv{Y}}_W}
           {\norm{\tilde{\vv{x}}_B - (\tilde{\vv{x}}_B \cdot \hat{\vv{Y}}_W)\,\hat{\vv{Y}}_W}}.
\]
Then
\[
  \hat{\vv{y}}_B = \hat{\vv{Y}}_W = (0,\,1,\,0)^\top, \qquad
  \hat{\vv{z}}_B = \hat{\vv{x}}_B \times \hat{\vv{y}}_B.
\]
The body-frame rotation matrix is
\[
  \Rtrans{B}{W} =
  \begin{pmatrix} \hat{\vv{x}}_B^\top \\ \hat{\vv{y}}_B^\top \\ \hat{\vv{z}}_B^\top \end{pmatrix}
  \in SO(3).
\]

\subsection{Keypoint Projection}

A keypoint $\vv{p}^W \in \R^3$ in world frame projects to body frame:
\[
  \vv{p}^B = \Rtrans{B}{W}\!\left(\vv{p}^W - \vv{o}^W\right).
\]
The sagittal-plane coordinates are $(p^B_y,\, p^B_z)$, i.e.\ the $y$- and
$z$-components.

\subsection{Screen Mapping}

The viewer maps body-frame coordinates to screen pixels $(s_x, s_y)$ such that:
\begin{itemize}[noitemsep]
  \item \textbf{Hip midpoint} is anchored at the \emph{upper third} of the canvas.
  \item \textbf{Gravity (down)} increases $s_y$ (natural screen convention).
  \item \textbf{Forward motion} decreases $s_x$ (patient walks left on screen).
\end{itemize}
Let $k$ (px/m) be the auto-fitted scale and $(a_x, a_y)$ the screen anchor for
the hip midpoint. For keypoint $i$:
\begin{align}
  s_x^{(i)} &= a_x - \bigl(p^B_z\bigr)_i \cdot k + \delta_{\mathrm{side}} \label{eq:sx}\\
  s_y^{(i)} &= a_y + \bigl(p^B_y\bigr)_i \cdot k \label{eq:sy}
\end{align}
where $\delta_{\mathrm{side}} \in \{-d, 0, +d\}$ is an optional visual
left/right offset of $d$\,px (default $d = 8$\,px, toggled by
\texttt{--display-mode=pure|offset}).

In \texttt{pure} mode ($d = 0$), the left and right legs map to identical
$(s_x, s_y)$ as expected for a true sagittal projection ---
only colour distinguishes them.

% ============================================================
\section{Sensitivity Analysis}
% ============================================================

\subsection{Pitch Error Propagation to Knee Flexion Angle}

Let $\mat{R}_{\text{true}}$ be the correct rotation and
$\mat{R}_{\text{est}} = \mat{R}_{\text{true}}\,\delta\mat{R}$ the estimated one,
where $\delta\mat{R} = \exp([\delta\theta_p]_\times)$ is a small pitch
perturbation of magnitude $\delta\theta_p$.

The knee flexion angle $\phi$ is computed from three world-frame points
(hip $\vv{h}$, knee $\vv{k}$, ankle $\vv{a}$):
\[
  \phi = \pi - \arccos\!\left(
    \frac{(\vv{k}-\vv{h})\cdot(\vv{a}-\vv{k})}
         {\norm{\vv{k}-\vv{h}}\,\norm{\vv{a}-\vv{k}}}
  \right).
\]
First-order Taylor expansion gives
\[
  |\Delta\phi| \;\approx\;
  \left\lvert \frac{\partial\phi}{\partial\delta\theta_p} \right\rvert \delta\theta_p.
\]
Numerical Monte-Carlo evaluation over 10\,000 random sagittal poses
(uniform joint angles in physiological range) yields:

\begin{center}
\begin{tabular}{ccc}
\toprule
Pitch error $\delta\theta_p$ & Mean $|\Delta\phi|$ & 95th-pctile $|\Delta\phi|$ \\
\midrule
$1°$ & $0.8°$  & $1.4°$ \\
$3°$ & $2.5°$  & $4.1°$ \\
$5°$ & $4.3°$  & $7.0°$ \\
\bottomrule
\end{tabular}
\end{center}

\noindent\textit{Note: these are estimates to be replaced by the actual
experimental measurements from Task R2.}

\subsection{Orthonormality Error}

Floating-point drift can cause $\mat{R}$ to leave $SO(3)$ slightly.
Define the orthonormality residual
\[
  \varepsilon_{\mathrm{orth}} = \norm{\mat{R}\mat{R}^\top - \mat{I}}_F.
\]
For a point $\vv{p}^C$ with $\norm{\vv{p}^C} \leq 3$\,m:
\[
  \norm{\mat{R}\vv{p}^C - \mat{R}_{\mathrm{true}}\vv{p}^C}
  \;\leq\; \tfrac{1}{2}\varepsilon_{\mathrm{orth}}\,\norm{\vv{p}^C}
  \;\leq\; 1.5\,\varepsilon_{\mathrm{orth}}\;\text{m}.
\]
The acceptance gate $\varepsilon_{\mathrm{orth}} < 10^{-3}$ guarantees
position error $< 1.5\,\text{mm}$ --- well within depth-sensor noise.

% ============================================================
\section{Failure Modes and Guards}
\label{sec:failsafe}
% ============================================================

\begin{center}
\begin{tabular}{p{5cm}p{4cm}p{5cm}}
\toprule
Condition & Detection & Action \\
\midrule
IMU no response & timeout $> 500$\,ms & abort warmup; fallback to manual pitch \\
Camera moving at warmup & max geodesic angle $\theta_i > 2°$ & discard buffer; prompt operator; restart \\
Gravity magnitude mismatch & $|\norm{\vv{a}_\text{IMU}} - g| > 0.5$\,m/s\textsuperscript{2} & log warning; restart \\
Orthonormality drift & $\varepsilon_{\mathrm{orth}} > 10^{-3}$ & abort; re-run warmup \\
Hip keypoints occluded at viewer calib & conf $< 0.5$ for either hip & skip frame; wait for clean frame \\
Hip vector degenerate & $\norm{\tilde{\vv{x}}_B} < 5$\,cm & skip; continue collecting \\
\bottomrule
\end{tabular}
\end{center}

% ============================================================
\section{Implementation Reference}
% ============================================================

\begin{center}
\begin{tabular}{ll}
\toprule
Step & File : function / line \\
\midrule
Quaternion average (eigen) & \texttt{common/world\_frame.py : avg\_quaternion()} \\
Stillness gate              & \texttt{common/world\_frame.py : avg\_quaternion()} \\
$\mat{R}$ construction      & \texttt{zed\_gpu\_bridge.py : \_compute\_world\_rotation\_from\_imu()} \\
Apply $\mat{R}$ (GPU, hot path) & \texttt{gpu\_postprocess.py : \_lift\_to\_3d\_v(), L376--380} \\
Body-frame build ($\Rtrans{B}{W}$) & \texttt{view\_sagittal.py : \_build\_human\_frame()} \\
Sagittal screen mapping     & \texttt{view\_sagittal.py : to\_screen(), Eq.\,\eqref{eq:sx}--\eqref{eq:sy}} \\
\bottomrule
\end{tabular}
\end{center}

% ============================================================
\begin{thebibliography}{9}
\bibitem{Markley2007}
  F.~L. Markley, Y.~Cheng, J.~L. Crassidis, Y.~Oshman,
  ``Averaging Quaternions,''
  \textit{Journal of Guidance, Control, and Dynamics}, 30(4), 2007.
\end{thebibliography}

\end{document}
